	% ============================================================
%  rk-test.tex
%  Test suite for the rkdiagram LaTeX package
%  Tests are organized by structural complexity.
%  Compile with: lualatex rk-test.tex
% ============================================================

\documentclass[12pt]{article}

% ---- Package under test ------------------------------------
\usepackage[
  spacing   = 1.3cm,
  modDepth  = 1.0cm,
  font      = \rmfamily,
  linewidth = 0.5pt,
  color     = black
]{rkdiagram}

% ---- Support packages --------------------------------------
\usepackage[letterpaper]{geometry}
\geometry{margin=1in}
\usepackage{graphicx}
\usepackage{booktabs}
\usepackage{microtype}
\usepackage{parskip}
\usepackage{xcolor}

% ---- Document styling --------------------------------------
\newcommand{\testsection}[1]{%
  \bigskip
  \noindent\rule{\linewidth}{0.8pt}\par\vspace{2pt}%
  {\noindent\raggedright\large\bfseries #1\par}\vspace{4pt}%
}

\newcommand{\testcase}[2]{%
  \medskip
  {\noindent\raggedright\textbf{Test \thesubsection:} \textit{``#1''}\par}%
  \stepcounter{subsection}\smallskip
  #2
  \par\medskip
}

% Manual counter for test numbering (subsection already defined by article class)
\setcounter{subsection}{1}

% ============================================================
\begin{document}

\begin{center}
  {\LARGE\bfseries \texttt{rkdiagram} Package Test Suite}\\[6pt]
  {\large Version 1.0 $\cdot$ Reed-Kellogg Sentence Diagramming}\\[4pt]
  {\small Compile with \texttt{lualatex}. Each diagram is followed by
   a notes line describing what structural feature it exercises.}
\end{center}

\bigskip

% ============================================================
\testsection{Group 1 — Baseline Only (Subject + Verb)}
% ============================================================
% These are the simplest possible diagrams.
% Expected output: one horizontal baseline, one full vertical divider.

\testcase{Dogs bark.}{%
  \begin{rkdiagram}
    \rkSubject{Dogs}
    \rkVerb{bark}
  \end{rkdiagram}
  \par\textit{Tests: bare subject, intransitive verb, no modifiers.}
}

\testcase{She laughed.}{%
  \begin{rkdiagram}
    \rkSubject{She}
    \rkVerb{laughed}
  \end{rkdiagram}
  \par\textit{Tests: pronoun subject, simple past verb.}
}

% ============================================================
\testsection{Group 2 — Modifiers on Subject and Verb}
% ============================================================
% Adjectives hang on diagonals below nouns.
% Adverbs hang on diagonals below verbs (or adjectives/adverbs).

\testcase{The old dog barked.}{%
  \begin{rkdiagram}
    \rkSubject[
      \rkAdj{the},
      \rkAdj{old}
    ]{dog}
    \rkVerb{barked}
  \end{rkdiagram}
  \par\textit{Tests: article + adjective both modifying same noun.}
}

\testcase{The dog barked loudly.}{%
  \begin{rkdiagram}
    \rkSubject[
      \rkAdj{The}
    ]{dog}
    \rkVerb[
      \rkAdv{loudly}
    ]{barked}
  \end{rkdiagram}
  \par\textit{Tests: adverb modifying verb.}
}

\testcase{The very old dog barked quite loudly.}{%
  \begin{rkdiagram}
    \rkSubject[
      \rkAdj{The},
      \rkAdj[
        \rkAdv{very}
      ]{old}
    ]{dog}
    \rkVerb[
      \rkAdv[
        \rkAdv{quite}
      ]{loudly}
    ]{barked}
  \end{rkdiagram}
  \par\textit{Tests: nested modifiers — adverb modifying adjective; adverb modifying adverb.}
}

% ============================================================
\testsection{Group 3 — Transitive Verbs (Direct Object)}
% ============================================================
% Direct object sits right of a SHORT vertical divider after the verb.

\testcase{The cat chased the mouse.}{%
  \begin{rkdiagram}
    \rkSubject[
      \rkAdj{The}
    ]{cat}
    \rkVerb{chased}
    \rkDObj[
      \rkAdj{the}
    ]{mouse}
  \end{rkdiagram}
  \par\textit{Tests: direct object with article modifier.}
}

\testcase{Maria writes beautiful letters.}{%
  \begin{rkdiagram}
    \rkSubject{Maria}
    \rkVerb{writes}
    \rkDObj[
      \rkAdj{beautiful}
    ]{letters}
  \end{rkdiagram}
  \par\textit{Tests: adjective on direct object.}
}

% ============================================================
\testsection{Group 4 — Indirect Object}
% ============================================================
% Indirect object sits on an angled platform below the verb.

\testcase{She gave him a gift.}{%
  \begin{rkdiagram}
    \rkSubject{She}
    \rkVerb{gave}
    \rkIObj{him}
    \rkDObj[
      \rkAdj{a}
    ]{gift}
  \end{rkdiagram}
  \par\textit{Tests: indirect object (pronoun) + direct object.}
}

\testcase{The teacher sent the students a warning.}{%
  \begin{rkdiagram}
    \rkSubject[
      \rkAdj{The}
    ]{teacher}
    \rkVerb{sent}
    \rkIObj[
      \rkAdj{the}
    ]{students}
    \rkDObj[
      \rkAdj{a}
    ]{warning}
  \end{rkdiagram}
  \par\textit{Tests: indirect object with article modifier.}
}

% ============================================================
\testsection{Group 5 — Linking Verbs}
% ============================================================
% Predicate nominatives and adjectives use a forward-diagonal (/)
% divider after the linking verb.

\testcase{The soup smells delicious.}{%
  \begin{rkdiagram}
    \rkSubject[
      \rkAdj{The}
    ]{soup}
    \rkVerb{smells}
    \rkPredAdj{delicious}
  \end{rkdiagram}
  \par\textit{Tests: predicate adjective after linking verb.}
}

\testcase{He is a teacher.}{%
  \begin{rkdiagram}
    \rkSubject{He}
    \rkVerb{is}
    \rkPredNom[
      \rkAdj{a}
    ]{teacher}
  \end{rkdiagram}
  \par\textit{Tests: predicate nominative with article.}
}

\testcase{The results seem surprisingly accurate.}{%
  \begin{rkdiagram}
    \rkSubject[
      \rkAdj{The}
    ]{results}
    \rkVerb{seem}
    \rkPredAdj[
      \rkAdv{surprisingly}
    ]{accurate}
  \end{rkdiagram}
  \par\textit{Tests: predicate adjective modified by adverb.}
}

% ============================================================
\testsection{Group 6 — Prepositional Phrases}
% ============================================================
% A slant line descends from the modified word; preposition sits
% ON the slant; object sits on a horizontal extension.

\testcase{The cat sat on the mat.}{%
  \begin{rkdiagram}
    \rkSubject[
      \rkAdj{The}
    ]{cat}
    \rkVerb[
      \rkPrep[
        \rkAdj{the}
      ]{on}{mat}
    ]{sat}
  \end{rkdiagram}
  \par\textit{Tests: prepositional phrase modifying verb.}
}

\testcase{The bird in the tree sang.}{%
  \begin{rkdiagram}
    \rkSubject[
      \rkAdj{The},
      \rkPrep[
        \rkAdj{the}
      ]{in}{tree}
    ]{bird}
    \rkVerb{sang}
  \end{rkdiagram}
  \par\textit{Tests: prepositional phrase modifying subject noun.}
}

\testcase{She ran into the old barn near the river.}{%
  \begin{rkdiagram}
    \rkSubject{She}
    \rkVerb[
      \rkPrep[
        \rkAdj{the},
        \rkAdj{old},
        \rkPrep[
          \rkAdj{the}
        ]{near}{river}
      ]{into}{barn}
    ]{ran}
  \end{rkdiagram}
  \par\textit{Tests: prepositional phrase modifying the object of another prepositional phrase.}
}

\testcase{The book on the shelf in the corner is mine.}{%
  \begin{rkdiagram}
    \rkSubject[
      \rkAdj{The},
      \rkPrep[
        \rkAdj{the},
        \rkPrep[
          \rkAdj{the}
        ]{in}{corner}
      ]{on}{shelf}
    ]{book}
    \rkVerb{is}
    \rkPredNom{mine}
  \end{rkdiagram}
  \par\textit{Tests: nested prepositional phrases (prep phrase modifying object of another prep).}
}

% ============================================================
\testsection{Group 7 — Compound Elements}
% ============================================================
% Compound subjects/verbs/objects use a dashed fork with the
% conjunction label on a dotted crossbar.

\testcase{John and Mary sang.}{%
  \begin{rkdiagram}
    \rkSubject{
      \rkCompound[and]{John}{Mary}
    }
    \rkVerb{sang}
  \end{rkdiagram}
  \par\textit{Tests: compound subject with coordinating conjunction.}
}

\testcase{She sang and danced.}{%
  \begin{rkdiagram}
    \rkSubject{She}
    \rkVerb{
      \rkCompound[and]{sang}{danced}
    }
  \end{rkdiagram}
  \par\textit{Tests: compound predicate.}
}

\testcase{The dog chased cats and birds.}{%
  \begin{rkdiagram}
    \rkSubject[
      \rkAdj{The}
    ]{dog}
    \rkVerb{chased}
    \rkDObj{
      \rkCompound[and]{cats}{birds}
    }
  \end{rkdiagram}
  \par\textit{Tests: compound direct object.}
}

\testcase{He painted the wall but not the ceiling.}{%
  \begin{rkdiagram}
    \rkSubject{He}
    \rkVerb{painted}
    \rkDObj{%
      \rkCompound[but not][
        \rkAdj{the}
      ]{wall}[
        \rkAdj{the}
      ]{ceiling}%
    }
  \end{rkdiagram}
  \par\textit{Tests: compound object with two-word conjunction \textit{but not}.}
}

% ============================================================
\testsection{Group 8 — Appositives}
% ============================================================
% Appositives appear in parentheses beside the word they rename.

\testcase{My brother, the doctor, arrived.}{%
  \begin{rkdiagram}
    \rkSubject[
      \rkAdj{My},
      \rkAppos{
        \rkAdj{the}doctor
      }
    ]{brother}
    \rkVerb{arrived}
  \end{rkdiagram}
  \par\textit{Tests: appositive renaming the subject.}
}

\testcase{We visited Paris, the city of light.}{%
  \begin{rkdiagram}
    \rkSubject{We}
    \rkVerb{visited}
    \rkDObj[
      \rkAppos{
        \rkAdj{the}city\rkPrep[
          \rkAdj{}
        ]{of}{light}
      }
    ]{Paris}
  \end{rkdiagram}
  \par\textit{Tests: appositive (with embedded prep phrase) on direct object.}
}

% ============================================================
\testsection{Group 9 — Subordinate Clauses}
% ============================================================
% A subordinate clause gets its own mini-baseline below the main
% one, connected by a dashed vertical line.

\testcase{The man who wore the hat left.}{%
  \begin{rkdiagram}
    \rkSubject[
      \rkAdj{The},
      \rkSubClause[who]{%
        \rkVerb{wore}%
        \rkDObj[
          \rkAdj{the}
        ]{hat}%
      }%
    ]{man}
    \rkVerb{left}
  \end{rkdiagram}
  \par\textit{Tests: relative clause modifying subject noun.}
}

\testcase{She knows that he lied.}{%
  \begin{rkdiagram}
    \rkSubject{She}
    \rkVerb{knows}
    \rkDObj{%
      \rkSubClause[that]{%
        \rkSubject{he}%
        \rkVerb{lied}%
      }%
    }
  \end{rkdiagram}
  \par\textit{Tests: noun clause functioning as direct object.}
}

\testcase{Because the rain fell, we stayed inside.}{%
  \begin{rkdiagram}
    \rkSubject{we}
    \rkVerb[
      \rkAdv{inside},
      \rkSubClause[because]{%
        \rkSubject[
          \rkAdj{the}
        ]{rain}%
        \rkVerb{fell}%
      }%
    ]{stayed}
  \end{rkdiagram}
  \par\textit{Tests: adverbial clause (fronted) modifying main verb.}
}

% ============================================================
\testsection{Group 10 — Complex Combined Structures}
% ============================================================
% These sentences exercise multiple features simultaneously.

\testcase{The tall woman in the red coat gave her young daughter
          a very old book.}{%
  \begin{rkdiagram}
    \rkSubject[
      \rkAdj{The},
      \rkAdj{tall},
      \rkPrep[
        \rkAdj{the},
        \rkAdj{red}
      ]{in}{coat}
    ]{woman}
    \rkVerb{gave}
    \rkIObj[
      \rkAdj{her},
      \rkAdj{young}
    ]{daughter}
    \rkDObj[
      \rkAdj{a},
      \rkAdj[
        \rkAdv{very}
      ]{old}
    ]{book}
  \end{rkdiagram}
  \par\textit{Tests: multiple adjectives, prep phrase on subject, indirect object
  with modifiers, adverb-modified adjective on direct object.}
}

\testcase{John and Maria quickly found the lost dog and gently
          carried it home.}{%
  \begin{rkdiagram}
    \rkSubject{
      \rkCompound[and]{John}{Maria}
    }
    \rkVerb{%
      \rkCompound[and][
        \rkAdv{quickly},
        \rkDObj[
          \rkAdj{the},
          \rkAdj{lost}
        ]{dog}
      ]{found}[
        \rkAdv{gently},
        \rkDObj{it},
        \rkAdv{home}
      ]{carried}%
    }
  \end{rkdiagram}
  \par\textit{Tests: compound subject + compound predicate, each verb with its
  own adverb, multiple direct objects.}
}

% ============================================================
\testsection{Group 11 — Edge Cases and Stress Tests}
% ============================================================

\testcase{Go!}{%
  \begin{rkdiagram}
    \rkSubject{(you)}   % implied subject — author inserts it in parens
    \rkVerb{Go}
  \end{rkdiagram}
  \par\textit{Tests: imperative with implied subject in parentheses (author convention).}
}

\testcase{Birds fly.}{%
  \begin{rkdiagram}
    \rkSubject{Birds}
    \rkVerb{fly}
  \end{rkdiagram}
  \par\textit{Tests: bare minimum diagram — no modifiers, no objects. Baseline
  and divider only.}
}

\testcase{New York City never sleeps.}{%
  \begin{rkdiagram}
    \rkSubject{New York City}
    \rkVerb[
      \rkAdv{never}
    ]{sleeps}
  \end{rkdiagram}
  \par\textit{Tests: multi-word proper noun in subject slot; adverb on verb.}
}

\testcase{She has been running every morning.}{%
  \begin{rkdiagram}
    \rkSubject{She}
    \rkVerb[
      \rkPrep[
        \rkAdj{every}
      ]{---}{morning}
    ]{has been running}
  \end{rkdiagram}
  \par\textit{Tests: multi-word verb phrase (auxiliary + participle) in single
  \texttt{\textbackslash rkVerb} slot; temporal prep phrase.}
}

\testcase{He never tells the truth.}{%
  \begin{rkdiagram}
    \rkSubject{He}
    \rkVerb[
      \rkAdv{never}
    ]{tells}
    \rkDObj[
      \rkAdj{the}
    ]{truth}
  \end{rkdiagram}
  \par\textit{Tests: negative adverb modifying verb.}
}

% ============================================================
\testsection{Group 12 — Gerunds}
% ============================================================
% A gerund sits on a stepped T-platform: word on a short top line,
% connected by a vertical post to a longer foot that rests on the
% noun-slot baseline.

\testcase{Swimming is fun.}{%
  \begin{rkdiagram}
    \rkSubject{
      \rkGerund{Swimming}
    }
    \rkVerb{is}
    \rkPredAdj{fun}
  \end{rkdiagram}
  \par\textit{Tests: gerund as subject.}
}

\testcase{She enjoys swimming.}{%
  \begin{rkdiagram}
    \rkSubject{She}
    \rkVerb{enjoys}
    \rkDObj{
      \rkGerund{swimming}
    }
  \end{rkdiagram}
  \par\textit{Tests: gerund as direct object.}
}

\testcase{His hobby is collecting stamps.}{%
  \begin{rkdiagram}
    \rkSubject[
      \rkAdj{His}
    ]{hobby}
    \rkVerb{is}
    \rkPredNom{
      \rkGerund{collecting}
    }
  \end{rkdiagram}
  \par\textit{Tests: gerund as predicate nominative.}
}

\testcase{Rapid swimming exhausts him.}{%
  \begin{rkdiagram}
    \rkSubject{
      \rkGerund[
        \rkAdj{Rapid}
      ]{swimming}
    }
    \rkVerb{exhausts}
    \rkDObj{him}
  \end{rkdiagram}
  \par\textit{Tests: gerund as subject with adjective modifier (post auto-height).}
}

\testcase{She is good at running quickly.}{%
  \begin{rkdiagram}
    \rkSubject{She}
    \rkVerb{is}
    \rkPredAdj[
      \rkPrep{at}{
        \rkGerund[
          \rkAdv{quickly}
        ]{running}
      }
    ]{good}
  \end{rkdiagram}
  \par\textit{Tests: gerund (with adverb modifier) as object of a preposition.}
}

\testcase{Swimming is fun. (forced post height via gerundPostH)}{%
  \begin{rkdiagram}
    \rkset{gerundPostH=2.5cm}
    \rkSubject{
      \rkGerund{Swimming}
    }
    \rkVerb{is}
    \rkPredAdj{fun}
  \end{rkdiagram}
  \par\textit{Tests: \texttt{\textbackslash rkset\{gerundPostH=2.5cm\}} forces the vertical bar
    to 2.5\,cm regardless of word height or modifiers.  Reset with
    \texttt{gerundPostH=0pt} to restore adaptive behaviour.}
}

\testcase{Rapid swimming exhausts him. (extended lower shelf via gerundShelfExt)}{%
  \begin{rkdiagram}
    \rkset{gerundShelfExt=1.5cm}
    \rkSubject{
      \rkGerund[
        \rkAdj{Rapid}
      ]{swimming}
    }
    \rkVerb{exhausts}
    \rkDObj{him}
  \end{rkdiagram}
  \par\textit{Tests: \texttt{\textbackslash rkset\{gerundShelfExt=1.5cm\}} extends the lower
    shelf and triangle foot 1.5\,cm to the right, giving modifier slants more room.}
}

\testcase{Rapid swimming exhausts him. (both gerundPostH and gerundShelfExt combined)}{%
  \begin{rkdiagram}
    \rkset{gerundPostH=2.0cm, gerundShelfExt=1.0cm}
    \rkSubject{
      \rkGerund[
        \rkAdj{Rapid}
      ]{swimming}
    }
    \rkVerb{exhausts}
    \rkDObj{him}
  \end{rkdiagram}
  \par\textit{Tests: both gerund controls combined.  Reset both with
    \texttt{\textbackslash rkset\{gerundPostH=0pt, gerundShelfExt=0pt\}} afterwards.}
}

% Reset gerund overrides after Group 12 tests
\rkset{gerundPostH=0pt, gerundShelfExt=0pt}

% ============================================================
\testsection{Group 13 — Participles}
% ============================================================
% A participle functions as an adjective modifying a noun.
% It sits on a hockey-stick shape: diagonal slanting down-right
% (word on it) then a short horizontal tail at the elbow.
% The tail can carry a participial object and/or modifiers.

\testcase{The barking dog scared us.}{%
  \begin{rkdiagram}
    \rkSubject[
      \rkAdj{The},
      \rkParticiple{barking}
    ]{dog}
    \rkVerb{scared}
    \rkDObj{us}
  \end{rkdiagram}
  \par\textit{Tests: bare participle (no tail) modifying subject noun.}
}

\testcase{The man running the race smiled.}{%
  \begin{rkdiagram}
    \rkSubject[
      \rkAdj{The},
      \rkParticiple[
        \rkPartObj[
          \rkAdj{the}
        ]{race}
      ]{running}
    ]{man}
    \rkVerb{smiled}
  \end{rkdiagram}
  \par\textit{Tests: participle with participial object on horizontal tail.}
}

\testcase{Exhausted by the climb, she rested.}{%
  \begin{rkdiagram}
    \rkSubject[
      \rkParticiple[
        \rkPrep[
          \rkAdj{the}
        ]{by}{climb}
      ]{Exhausted}
    ]{she}
    \rkVerb{rested}
  \end{rkdiagram}
  \par\textit{Tests: participle with a prepositional phrase on the tail.}
}

\testcase{The child sleeping quietly in the corner dreamed.}{%
  \begin{rkdiagram}
    \rkSubject[
      \rkAdj{The},
      \rkParticiple[
        \rkAdv{quietly},
        \rkPrep[
          \rkAdj{the}
        ]{in}{corner}
      ]{sleeping}
    ]{child}
    \rkVerb{dreamed}
  \end{rkdiagram}
  \par\textit{Tests: participle with adverb and prepositional phrase on the tail.}
}

\testcase{The dog chased the cat sitting on the fence.}{%
  \begin{rkdiagram}
    \rkSubject[
      \rkAdj{The}
    ]{dog}
    \rkVerb{chased}
    \rkDObj[
      \rkAdj{the},
      \rkParticiple[
        \rkPrep[
          \rkAdj{the}
        ]{on}{fence}
      ]{sitting}
    ]{cat}
  \end{rkdiagram}
  \par\textit{Tests: participle modifying the direct object, with prep phrase on tail.}
}

% ============================================================
\testsection{Group 14 — Infinitives}
% ============================================================

\testcase{To run is fun.}{%
  \begin{rkdiagram}
    \rkSubject{
      \rkInfinitive{run}
    }
    \rkVerb{is}
    \rkPredNom{fun}
  \end{rkdiagram}
  \par\textit{Tests: infinitive as subject, no mods.}
}

\testcase{She wants to win.}{%
  \begin{rkdiagram}
    \rkSubject{She}
    \rkVerb{wants}
    \rkDObj{
      \rkInfinitive{win}
    }
  \end{rkdiagram}
  \par\textit{Tests: infinitive as direct object, no mods.}
}

\testcase{To run quickly is fun.}{%
  \begin{rkdiagram}
    \rkSubject{
      \rkInfinitive[
        \rkAdv{quickly}
      ]{run}
    }
    \rkVerb{is}
    \rkPredNom{fun}
  \end{rkdiagram}
  \par\textit{Tests: infinitive as subject with adverb modifier (adaptive post height).}
}

\testcase{He wants to run in the race.}{%
  \begin{rkdiagram}
    \rkSubject{He}
    \rkVerb{wants}
    \rkDObj{
      \rkInfinitive[
        \rkPrep[
          \rkAdj{the}
        ]{in}{race}
      ]{run}
    }
  \end{rkdiagram}
  \par\textit{Tests: infinitive as direct object with prep phrase modifier on verb.}
}

\testcase{She came to win a prize.}{%
  \begin{rkdiagram}
    \rkSubject{She}
    \rkVerb[
      \rkInfinitive{win}
    ]{came}
    \rkDObj{prize}
  \end{rkdiagram}
  \par\textit{Tests: infinitive as adverbial modifier on verb, no tail.}
}

\testcase{He found books to read.}{%
  \begin{rkdiagram}
    \rkSubject{He}
    \rkVerb{found}
    \rkDObj[
      \rkInfinitive{read}
    ]{books}
  \end{rkdiagram}
  \par\textit{Tests: infinitive as adjectival modifier on direct object.}
}

\testcase{She came to win the prize.}{%
  \begin{rkdiagram}
    \rkSubject{She}
    \rkVerb[
      \rkInfinitive[
        \rkInfObj{prize}
      ]{win}
    ]{came}
  \end{rkdiagram}
  \par\textit{Tests: infinitive modifier with direct object on mini-baseline.}
}

\testcase{He trained to finish the race quickly.}{%
  \begin{rkdiagram}
    \rkSubject{He}
    \rkVerb[
      \rkInfinitive[
        \rkInfObj[
          \rkAdj{the}
        ]{race}
        \rkAdv{quickly}
      ]{finish}
    ]{trained}
  \end{rkdiagram}
  \par\textit{Tests: infinitive modifier with object and adverb on mini-baseline.}
}

% ============================================================
\testsection{Group 15 — Nominative Absolutes}
% ============================================================

\testcase{The battle over, the soldiers retreated.}{%
  \begin{rkdiagram}
    \rkSubject[
      \rkAdj{The}
    ]{soldiers}
    \rkVerb{retreated}
    \rkNomAbs[
      \rkAdj{The}
    ]{battle}{}{over}
  \end{rkdiagram}
  \par\textit{Tests: bare nominative absolute with adjective on noun.}
}

\testcase{The sun setting, they departed.}{%
  \begin{rkdiagram}
    \rkSubject{they}
    \rkVerb{departed}
    \rkNomAbs[
      \rkAdj{The}
    ]{sun}{}{setting}
  \end{rkdiagram}
  \par\textit{Tests: nominative absolute with adjective on noun.}
}

\testcase{His work finished, he ran.}{%
  \begin{rkdiagram}
    \rkSubject{he}
    \rkVerb{ran}
    \rkNomAbs[
      \rkAdj{His}
    ]{work}{\rkPartObj{done}}{finished}
  \end{rkdiagram}
  \par\textit{Tests: NA with participial object on tail.}
}

\testcase{His staff grasped in his hand, the old man walked the long road slowly.}{%
  \begin{rkdiagram}
    \rkSubject[
      \rkAdj{The}
      \rkAdj{old}
    ]{man}
    \rkVerb[
      \rkAdv{slowly}
    ]{walked}
    \rkDObj[
      \rkAdj{the}
      \rkAdj{long}
    ]{road}
    \rkNomAbsHeight{4cm}
    \rkNomAbs[
      \rkAdj{his}
    ]{staff}{
      \rkPrep[
        \rkAdj{his}
      ]{in}{hand}
    }{grasped}
  \end{rkdiagram}
  \par\textit{Tests: NA with modifier chains on main clause (vertical clearance).}
}

\testcase{Her hope gone, she wept.}{%
  \begin{rkdiagram}
    \rkSubject{she}
    \rkVerb{wept}
    \rkNomAbs[
      \rkAdj{Her}
    ]{hope}{}{gone}
  \end{rkdiagram}
  \par\textit{Tests: NA with no participle (empty ptcl-word, implied being).}
}

% ============================================================
\testsection{Group 16 — Manual Spacing Controls}
% ============================================================
% \rkSlotWidth{len}  sets a minimum line width for the next slot only.
% \rkSkip{len}       inserts a gap in the baseline before the next slot.
% Both are one-shot — they reset automatically after use.

\testcase{Dogs bark. (default spacing for reference)}{%
  \begin{rkdiagram}
    \rkSubject{Dogs}
    \rkVerb{bark}
  \end{rkdiagram}
  \par\textit{Baseline: auto-width lines, no adjustments.}
}

\testcase{Dogs bark. (wider subject line)}{%
  \begin{rkdiagram}
    \rkSlotWidth{4cm}
    \rkSubject{Dogs}
    \rkVerb{bark}
  \end{rkdiagram}
  \par\textit{\texttt{\textbackslash rkSlotWidth\{4cm\}} before \texttt{\textbackslash rkSubject}
    forces the subject line to at least 4\,cm.}
}

\testcase{Dogs bark. (gap before verb)}{%
  \begin{rkdiagram}
    \rkSubject{Dogs}
    \rkSkip{1.5cm}
    \rkVerb{bark}
  \end{rkdiagram}
  \par\textit{\texttt{\textbackslash rkSkip\{1.5cm\}} before \texttt{\textbackslash rkVerb}
    inserts a 1.5\,cm gap in the baseline (extends the subject line rightward).}
}

\testcase{The cat chased the mouse. (widen verb and object lines)}{%
  \begin{rkdiagram}
    \rkSubject[
      \rkAdj{The}
    ]{cat}
    \rkSlotWidth{3.5cm}
    \rkVerb{chased}
    \rkSlotWidth{3.5cm}
    \rkDObj[
      \rkAdj{the}
    ]{mouse}
  \end{rkdiagram}
  \par\textit{Both verb and object lines forced to 3.5\,cm minimum — useful for
    aligning words on parallel diagrams or leaving room for annotations.}
}

\testcase{She gave him a gift. (skip before indirect object platform)}{%
  \begin{rkdiagram}
    \rkSubject{She}
    \rkSlotWidth{3cm}
    \rkVerb{gave}
    \rkIObj{him}
    \rkDObj[
      \rkAdj{a}
    ]{gift}
  \end{rkdiagram}
  \par\textit{\texttt{\textbackslash rkSlotWidth\{3cm\}} on the verb widens it
    to give the indirect-object platform more room.}
}

\testcase{The cat sat on the mat. (longer prep slant)}{%
  \begin{rkdiagram}
    \rkSubject[
      \rkAdj{The}
    ]{cat}
    \rkVerb[
      \rkSlantLen{2.5cm}
      \rkPrep[
        \rkAdj{the}
      ]{on}{mat}
    ]{sat}
  \end{rkdiagram}
  \par\textit{\texttt{\textbackslash rkSlantLen\{2.5cm\}} before \texttt{\textbackslash rkPrep}
    forces the ``on'' slant to 2.5\,cm — useful when the prep label is long or
    when more visual separation is needed.}
}

\testcase{Dogs bark. (combining skip and minwidth)}{%
  \begin{rkdiagram}
    \rkSlotWidth{3cm}
    \rkSubject{Dogs}
    \rkSkip{0.8cm}
    \rkSlotWidth{3cm}
    \rkVerb{bark}
  \end{rkdiagram}
  \par\textit{Both controls used together: 3\,cm subject, 0.8\,cm gap,
    3\,cm verb — each resets automatically after its slot.}
}

% ============================================================
\testsection{Group 17 — Mark Twain Stress Test}
% ============================================================

% Compound sentence (two independent clauses joined by ``and'').
% Diagrammed as two stacked baselines.
\testcase{A swarm of swarthy, noisy, lying, shoulder-shrugging, gesticulating
Portuguese boatmen, with brass rings in their ears and fraud in their hearts,
climbed the ship's sides, and various parties of us contracted with them to
take us ashore at so much a head, silver coin of any country.}{%

  \scalebox{0.6}{%
  \rkCompoundSent[and][10cm]%
    [spacing=1.0cm,modDepth=0.9cm,cmpSep=1.8cm]%
    {%
      \begin{rkdiagram}
        \rkSubject[
          \rkAdj{A},
          \rkset{cmpSep=1.0cm}
          \rkPrep[
            \rkAdj{swarthy},
            \rkAdj{noisy},
            \rkAdj{lying},
            \rkAdj{shoulder-shrugging},
            \rkAdj{gesticulating},
            \rkAdj{Portuguese},
            \rkPrep{with}{%
              \rkCompound[and][
                \rkAdj{brass},
                \rkPrep[
                  \rkAdj{their}
                ]{in}{ears}
              ]{rings}[
                \rkPrep[
                  \rkAdj{their}
                ]{in}{hearts}
              ]{fraud}%
            }%
          ]{of}{boatmen}%
        ]{swarm}
        \rkVerb{climbed}
        \rkDObj[
          \rkAdj{the},
          \rkAdj{ship's}
        ]{sides}
      \end{rkdiagram}%
    }%
    [spacing=1.0cm,modDepth=1.0cm]%
    {%
      \begin{rkdiagram}
        \rkSubject[
          \rkAdj{various},
          \rkPrep{of}{us}%
        ]{parties}
        \rkVerb[
          \rkPrep{with}{them},
          \rkSubClause[to]{%
            \rkVerb[
              \rkAdv{ashore},
              \rkPrep[
                \rkAdj{so~much},
                \rkAdj{a},
                \rkAppos{silver coin of any country}%
              ]{at}{head}%
            ]{take}
            \rkDObj{us}%
          }%
        ]{contracted}
      \end{rkdiagram}%
    }%
  }%
  \par\textit{(Mark Twain, \textit{The Innocents Abroad}, 1869)}
}

\end{document}
